\documentclass[11pt,a4paper]{article}
\usepackage[margin=25mm]{geometry}
\usepackage{amsmath,amssymb,amsthm,hyperref}
\newtheorem{lemma}{Lemma}
\newcommand{\gen}[1]{\langle#1\rangle}
\title{Saturated mixed block systems at a crossing}
\author{Human Compression project: bounded L044 proof component}
\date{}
\begin{document}\maketitle
Retain the unique-crossing setting in $G\leq A_n$, $n\geq12$,
with all maximal omissions intransitive. The two $G_i$-orbits
are $A,B$, with crossing endpoints $a,b$. The deficits and
minimal simultaneous block-image family $M$ are those of the
accepted block-deficit note. Block systems are chosen maximal,
so their induced actions are primitive. The arguments below
assume $\alpha=\beta=0$ unless stated otherwise.

\begin{lemma}[A binary and a maximal symmetric block image]
Suppose the block numbers are $2,m$, with $m\geq3$. The actual
$M$-action is faithfully the diagonal $C_2\times S_m$ action.
Its distinguished $C_2$ generator and its pure $S_m$ component
occur on opposite sides of $i$. The entire left and right
intervals of $G_i$ are supported on their respective original
orbits, after reversal if necessary.
\end{lemma}
\begin{proof}
Saturation gives the full block image $C_2\times S_m$ and makes
its lift faithful, by the paid actual-orbit argument. Choose a
minimal subfamily $T\subset M$ generating the $S_m$ quotient.
Normal-quotient normalization bounds the remaining independent
family by one. Whiston then forces $|T|=m-1$, leaving exactly
one distinguished generator $e$ outside $T$.

The maximal involutory string images of $T$ have a connected
noncommuting path. This follows from the accepted extremal
classification for $m\geq5$, and from the complete elementary
$S_3,S_4$ cases in the anchored symmetric lemma. Their indices
therefore form an interval. Since their image is transitive,
one of them moves the block containing $b$, so that interval
contains $i-1$ or $i+1$; reverse so it starts at $i+1$.
Only its first member moves that block.

The $S_m$ image of $e$ is trivial. If $e$ lies on the opposite
side of $i$, a nontrivial image would disconnect the active set
of this primitive block action, contrary to the paid primitive
interval lemma. If it lies on the same side, it is beyond $T$,
fixes the crossing block, and centralizes its far-end omission
(or all of $T$). The anchored symmetric lemma makes its image
trivial in this case too. Thus $e$ is the pure nonidentity
$C_2$ element. Its index must be $i-1$, since it exchanges
the two blocks containing $a$ and cannot lie beyond $T$.

The subgroup generated by $T$ intersects this central $C_2$
trivially, by IP in the faithful block action. It is a character
graph $S_m\to C_2$. That character is trivial: the $m-2$ tail
generators of $T$ fix the crossing block and are independent in
$S_{m-1}$, so generate it by Whiston's equality clause. Some
has odd $S_m$ image. All have indices at least $i+2$ and fix
$a$, so have trivial $C_2$ image. A nontrivial character of
$S_m$, namely its sign, is therefore impossible. The pure
components are on opposite sides as claimed.

Every generator on the $S_m$ side fixes $a$: its near generator
is now pure on the other original orbit, and the more distant
ones fix $a$ by the crossing property. The accepted commuting
anchor lemma makes this whole interval trivial on the orbit of
$a$. The symmetric statement holds for the other interval and
$b$. This proves exact support separation, not merely separation
of the block images.
\end{proof}

\begin{lemma}[The remaining saturated $D=2$ case]
The residual block systems $(k_A,m_A)=(4,3)$ and
$(k_B,m_B)=(k,2)$, with $\alpha=\beta=\gamma=0$, are impossible.
\end{lemma}
\begin{proof}
Lemma~1 separates the two interval supports. On the twelve-point
orbit, $M$ is the natural $S_3$ on four copies, with two
consecutive generators at the end next to $i$. Only its one
outer neighbor can fail to centralize all of this $S_3$.
Any further distinguished generator fixes the crossing endpoint
and therefore fixes its whole three-point copy. The full
centralizer on the remaining three copies is $S_3$ permuting
those copies. A transposition of three-point copies is odd on
the twelve points. Since this interval fixes the other original
orbit pointwise, its even part is $A_3$, which contains no
nonidentity involution. No such further generator exists.
This interval has rank at most three.

On the other orbit the last distinguished generator is its
pure swap of two $k$-point blocks. Apart from its possible
outer neighbor, every other generator centralizes that swap
and fixes the crossing endpoint; it acts faithfully on the
$k-1$ nonanchor row pairs. Whiston bounds their number by
$k-2$, so this interval has rank at most $k$. This includes
$k=2$. Thus the sum of the two interval ranks is at most $k+3$.
But the exact deficit formula under the stated saturation gives
$r=k+6$, hence that sum is $r-1=k+5$. This contradiction
excludes the case.
\end{proof}

\begin{lemma}[Saturated mixed two-block and two-point-block systems]
For block systems $(k,2)$ and $(2,m)$, $m\geq3$,
$\alpha=\beta=0$ implies $r\leq k+m-1=(n-2)/2$.
\end{lemma}
\begin{proof}
By Lemma~1 the original supports separate. On the two-point-block
orbit the $m-1$ generators of the pure $S_m$ component end
next to $i$. A farther generator centralizing all of $S_m$ fixes
the crossing endpoint and its two-point block pointwise.
Commutation transports this fixation to every block, so it is
trivial. The only possible additional generator is the one
outer neighbor $x$ of the $S_m$ interval. Its block image
centralizes the far-end omission and fixes the crossing block,
so is trivial by the anchored symmetric lemma. Pointwise
fixation of the two-point crossing block then propagates to
at least $m-1$ blocks. On at most one remaining two-point
block $x$ may be a transposition. Its action on the other
original orbit is trivial by support separation, so evenness
excludes that possibility too. Thus this interval has exactly
$m-1$ generators.

On the two-block orbit, the pure swap has $k$ transpositions,
so $k$ is even. Its outer neighbor must exist and must move
the opposite anchor to another row. If $k=2$ this is one
odd transposition, impossible. For $k\geq4$ the proof of the
accepted paired-row side-rank lemma applies without any
assumption on the number of blocks in the opposite orbit.
For clarity, its used hypotheses and endpoint step are as follows.
The diagonal tail has rank at most $k-2$, so the entire interval
has rank at most $k$. At equality, its tail is a maximal
$S_{k-1}$ string. For even $k\geq6$ the ordinary endpoint
omission makes its outer neighbor a single odd transposition.
For $k=4$ the complete two endpoint possibilities are
\[
 (1\ 2)(0'\ 3')\quad\hbox{and}\quad(1'\ 2')(0'\ 3').
\]
After deleting the far tail generator, the remaining near tail,
this endpoint, and the block swap are transitive on all eight
points: their contracted row edges are $2$--$3$, $0$--$3$,
and $1$--$2$. The entire opposite interval is unchanged and
transitive on its original orbit; $\rho_i$ joins them. The
global omission would be transitive, contradicting a standing
hypothesis. Thus the two-block interval has rank at most $k-1$.
Adding $\rho_i$ gives $r\leq(k-1)+(m-1)+1=k+m-1$.
\end{proof}

\begin{lemma}[The one-deficit mixed block image]
For block systems $(k,2)$ and $(2,m)$, $m\geq3$, the case
$\alpha=1$, $\beta=\gamma=0$ cannot violate the rank bound.
\end{lemma}
\begin{proof}
Suppose it does. The exact deficit gives $r=k+m$,
$|M|=m-1$, and $k$ distinguished generators outside $M$ and $i$.
The zero orbit deficit makes the block action faithful on $M$,
with diagonal natural actions on $k$ binary rows and two
$m$-point copies. Let $H$ be its primitive degree-$m$ block
image. Its active indices form one interval next to $i$ by
the accepted primitive interval lemma; reverse so that this
interval starts at $i+1$.

First suppose there is a member of $M$ outside that interval.
Its image on $H$ is trivial, so minimality makes it the sole
such member and the pure $C_2$ generator $\rho_{i-1}$.
The remaining $m-2$ generators project independently onto $H$:
any projected redundancy could be canceled by this available
pure binary generator. IP makes their subgroup a character
graph over $H$ in $C_2\times H$.
That graph is trivial on the binary coordinate. Otherwise
only its first generator could have nontrivial binary image,
since all its later generators fix the crossing endpoint.
Both neighbors of $i$ would then swap the paired anchor in
$A$, leaving that pair invariant under $G_i$ by the same
commutation argument used in the two-swap proof. This is
impossible. The commuting anchor lemma now separates the full
left and right supports.

The two-block side has rank at most $k-1$ by the complete
side-rank argument in Lemma~3 (including the exclusion at $k=2$).
The other side consists of its $m-2$ members of $M$ and at
most one outer neighbor. To justify the latter assertion,
a farther generator centralizes all of $H$ and fixes the
crossing endpoint. It therefore fixes the entire natural
copy containing that endpoint, and preserves the other copy.
An involution centralizing a primitive group of degree at least
three is the identity: if nontrivial it is fixed-point-free,
and its two-point orbits form an invariant nontrivial block
system. Thus the farther generator is also trivial on the
second copy. Support separation makes it trivial everywhere.
The total rank is consequently at most
$(k-1)+(m-1)+1=k+m-1$, a contradiction.

It remains that all of $M$ is the single interval starting
at $i+1$. Only its first generator has nontrivial binary
image, and only that generator can move the degree-$m$
crossing block. The accepted binary anchored-string lemma
applies: the $H$ image is $S_m$ and its far-end omission has
an orbit containing the crossing block of size at least
$m-1$, with centralizer fixing that block trivial. For
$m=3$, not stated in that lemma, the assertion is immediate:
the last generator is the transposition of the other two
blocks, the first projects to a transposition moving the
crossing block, and its one-generator far-end omission has
the required two-point orbit and centralizer property.

Let $u=\rho_{i-1}$ and $v=\rho_{i+m}$ when those indices
exist, and let $C$ be the other generators outside $M$ and $i$.
They centralize $M$ and fix both crossing endpoints. On the
two $m$-point copies their action is trivial: fixing one
point fixes the first copy, while the natural $S_m$ has
trivial centralizer on the second. On $A$ they are diagonal
and fix its anchor row. Hence their action is faithful in
$S_{k-1}$ and they number at most $k-2$.
There are $k$ outside generators, so both $u,v$ exist and
$C$ has exactly $k-2$ members. Whiston's equality clause gives
$\gen C=S_{k-1}$ on the nonanchor rows (trivial if $k=2$).

The generator $v$ fixes both crossing endpoints and centralizes
the far-end omission of $M$. Its degree-$m$ block image is
therefore trivial. It fixes the two-point crossing block
pointwise, and commutation propagates this fixation to at
least $m-1$ blocks. On the remaining orbit it can thus move
at most the final two-point block, by one transposition.
On $A$, the generator $v$ commutes with the first member
of $M$ (their index distance is $m-1\geq2$) and preserves
its two blocks. It acts diagonally and hence evenly there.
Evenness excludes a single transposition on the other
original orbit. Thus $v$ is trivial on that whole orbit.
Its remaining diagonal action on $A$ fixes the anchor row
and lies in the already full $\gen C=S_{k-1}$.
The distinguished generator $v$ is redundant, the final
contradiction.
\end{proof}

\paragraph{Boundary.}
These proofs retain the full accepted L047 independent-generation
and anchored symmetric interfaces, the L044 block-deficit proof,
and the two-block completion. They do not assert that a block
quotient has IP unless faithfulness has been obtained from
$\beta=0$. Only the mixed $\alpha=0,\beta=1,\gamma=0$ case remains
separate in this numerical residual family, as do the
other outstanding principal source obligations. Source framework:
Cameron--Fernandes--Leemans--Mixer, arXiv:1605.09173v1, Section~5.1.
No new census, Lean theorem, original or question credit is claimed.
\end{document}
